\section{Institutional Interpretation}\label{sec:institutional}

The model is intended for enterprise digital services in which users can make relationship-specific integrations before all future usage terms are known and in which both fixed-access and usage-linked charging are feasible. The architecture distinction $p=0$ versus $p>0$ is therefore an economic restriction on the marginal price of future usage, not a literal description of any one vendor's entire contract menu.

Several current market facts motivate this interpretation. OpenAI documents ChatGPT Business seats with fixed per-user charges and included usage limits, together with optional workspace credits that can extend eligible usage beyond those limits; Enterprise and Edu arrangements can instead draw from shared contract-level credit pools \citep{openaibusinesspricing2026,openaiflexible2026}. In cloud computing, AWS Savings Plans exchange discounted compute prices for a one- or three-year usage commitment \citep{awssavingsplans2026}, while Google Cloud committed-use discounts similarly exchange discounted pricing for minimum resource use or minimum spend commitments over one- or three-year horizons \citep{googlecud2026}. These examples establish coexistence of fixed-access, usage-sensitive, and commitment-based payment structures in closely related markets.

The timing in Section~\ref{sec:model} is a model interpretation rather than an empirical claim. An enterprise may connect an AI service to data, software, internal processes, or task-specific workflows before future pricing terms are fully settled. We summarize the current value of such integration by $b_j$ and its heterogeneous sunk cost by $k$. The model does not estimate those costs and does not infer from observed pricing menus that installed integration causally changes later tariff architecture.

The interpretation of $\mu$ is also deliberately narrow. It is a real cost of operating a metered architecture. It should not be reinterpreted as customer budget risk, procurement aversion, a transfer payment, or a reduced-form welfare penalty. Those objects would require different primitives.

Finally, the frozen baseline does not endogenize AI model quality. The parameters $a_L$ and $a_H$ summarize future task value. Any statement that model improvement moves the economy from flat pricing toward metering is therefore conditional on how an exogenous quality change moves the certified thresholds; no universal quality-improvement path is claimed.
